% Greek abstract: mirrors the English abstract (three paragraphs + keywords)

Οι εφαρμογές εκμάθησης γλωσσών με φορητές συσκευές (Mobile-Assisted Language Learning, MALL) αποτελούν σήμερα το πιο διαδεδομένο μέσο αυτόνομης γλωσσικής μάθησης, ωστόσο οι περισσότερες στηρίζονται σε ομοιόμορφες ασκήσεις αποκομμένες από το γλωσσικό πλαίσιο, «διακοσμημένες» με πόντους και σερί ημερών, ενώ η απώλεια χρηστών μέσα στην πρώτη εβδομάδα αναφέρεται ότι ξεπερνά το 60\%. Η παρούσα εργασία εξετάζει αν τα μέσα που οι μαθητές ήδη καταναλώνουν συστηματικά --- τα βίντεο μικρού μήκους και τα παιχνίδια hypercasual --- μπορούν να μετατραπούν σε ένα αποτελεσματικό, εξατομικευμένο περιβάλλον εκμάθησης λεξιλογίου. Παρουσιάζει το \textit{Glosy}, μια εφαρμογή για φορητές συσκευές πολλαπλών πλατφορμών (cross-platform), η οποία υποστηρίζει προς το παρόν Αγγλικά, Ελληνικά και Φαρσί σε οποιονδήποτε συνδυασμό μητρικής και εκμαθούμενης γλώσσας, με λεξικό περίπου 41.000 λέξεων επισημασμένων ανά επίπεδο επάρκειας.

Κεντρικό στοιχείο του σχεδιασμού είναι ένα μοντέλο μαθητή σε επίπεδο λέξης: κάθε λέξη που παρακολουθεί ο μαθητής φέρει μια κατάσταση μνήμης του Free Spaced Repetition Scheduler (FSRS) (δυσκολία, σταθερότητα και ανακτησιμότητα), η οποία ενημερώνεται στη συσκευή μετά από κάθε βαθμολογημένη επανάληψη, ώστε η εφαρμογή να λειτουργεί χωρίς σύνδεση και να συγχρονίζεται όταν αυτή επανέλθει. Οι μαθητές παρακολουθούν reels --- βίντεο διάρκειας έως 90 δευτερολέπτων με χρονισμένους, χωρισμένους σε λέξεις και μεταφρασμένους υπότιτλους, τα οποία μπορούν και οι ίδιοι να δημιουργήσουν --- και προσθέτουν λέξεις στο λεξιλόγιό τους πατώντας πάνω τους. Τα reels επιλέγονται από ένα σύστημα συστάσεων τριών σταδίων τύπου «καταρράκτη» (cascade). Ένα φίλτρο κατανοησιμότητας κρατά μόνο τα reels στα οποία ο μαθητής γνωρίζει ήδη τουλάχιστον το 98\% του λεξιλογίου· ένα στάδιο επανάληψης σε αραιά διαστήματα προτάσσει τα reels που περιέχουν λέξεις προς επανάληψη, ώστε η παρακολούθηση ενός reel να λειτουργεί ταυτόχρονα ως βαθμολογημένη επανάληψη· και ένα στάδιο βασισμένο στο περιεχόμενο ταξινομεί τα υπόλοιπα με βάση την ομοιότητα συνημιτόνου (TF-IDF) κάθε reel με ένα προφίλ που δημιουργείται από τις αλληλεπιδράσεις του μαθητή. Δύο παιχνίδια hypercasual, ένα πολύγλωσσο παιχνίδι τύπου Wordle και ένα διαδικαστικά παραγόμενο σταυρόλεξο με την ονομασία Word of Wonders, αντλούν τις λέξεις τους από το ίδιο λεξιλόγιο που παρακολουθεί ο μαθητής, ώστε κάθε γύρος να αποτελεί επίσης επανάληψη σε αραιά διαστήματα.

Δεν έχει ακόμη πραγματοποιηθεί ελεγχόμενη μελέτη με μαθητές, και για τον λόγο αυτό δεν υποστηρίζεται η παιδαγωγική αποτελεσματικότητα· αντί γι' αυτό, αξιολογούνται απευθείας τα αλγοριθμικά συστατικά του συστήματος. Σε 4.200 δοκιμές αξιολόγησης (benchmark) ο γεννήτορας του σταυρολέξου δεν απέτυχε ποτέ, και ο πλούτος των πλεγμάτων του περιοριζόταν από το όριο γραμμάτων και όχι από το όριο λέξεων. Το \textit{Glosy} αντιμετωπίζει επίσης το πρόβλημα ψυχρής εκκίνησης (cold start) μιας νέας εφαρμογής: χωρίς ακόμη βάση χρηστών, κανένας χρήστης δεν έχει ανεβάσει reels, οπότε η βιβλιοθήκη περιεχομένου είναι μικρή και η μηχανή συστάσεων δεν διαθέτει ιστορικό αλληλεπιδράσεων για εξατομίκευση. Η εργασία κλείνει με τους βασικούς περιορισμούς --- την απουσία μελέτης με χρήστες και την εκτός σύνδεσης επισημείωση των νέων reels --- και με κατευθύνσεις για μελλοντική εργασία: αυτοματοποιημένη επισημείωση, μετριασμό της ψυχρής εκκίνησης και τυχαιοποιημένη ελεγχόμενη αξιολόγηση.

\textbf{Λέξεις κλειδιά:} εκμάθηση γλωσσών με φορητές συσκευές, βίντεο μικρού μήκους, επανάληψη σε αραιά διαστήματα, συστήματα συστάσεων, παιχνίδια hypercasual, εκμάθηση λεξιλογίου
